\documentclass[12pt,a4paper]{amsart}

% Packages
\usepackage[utf8]{inputenc}
\usepackage[T1]{fontenc}
\usepackage{amsmath,amssymb,amsthm,amscd}
\usepackage{hyperref}
\usepackage{booktabs}
\usepackage{geometry}
\geometry{margin=2.5cm}

% Theorem environments
\theoremstyle{plain}
\newtheorem{theorem}{Theorem}[section]
\newtheorem{lemma}[theorem]{Lemma}
\newtheorem{proposition}[theorem]{Proposition}
\newtheorem{corollary}[theorem]{Corollary}

\theoremstyle{definition}
\newtheorem{definition}[theorem]{Definition}
\newtheorem{example}[theorem]{Example}
\newtheorem{remark}[theorem]{Remark}

% Custom commands
\newcommand{\Collatz}{\text{Collatz}}
\newcommand{\corrsum}{\operatorname{corrsum}}
\newcommand{\shift}{\operatorname{shift}}
\newcommand{\tail}{\operatorname{tail}}
\newcommand{\ord}{\operatorname{ord}}
\newcommand{\vp}{v_2}  % 2-adic valuation
\renewcommand{\gcd}{\operatorname{gcd}}
\newcommand{\lcm}{\operatorname{lcm}}
\renewcommand{\Im}{\operatorname{Im}}

\title{Cyclic Shift Obstructions for Non-Trivial Collatz Cycles}
\author{Eric Merle}
\date{March 2026}
\thanks{The author acknowledges discussions with members of the Lean formalization community.}

\begin{document}

\begin{abstract}
We establish the nonexistence of nontrivial periodic orbits of the Collatz map through two complementary approaches. The primary contribution is a rotation obstruction theorem: for any hypothetical cycle of length $k$, the correction sums of consecutive cyclic shifts satisfy a single algebraic identity forcing the divisibility condition $3^{2k} \equiv 1 \pmod{d(k)}$ where $d(k) = 2^S - 3^k$. Computational verification shows $3^{2k} \not\equiv 1 \pmod{d(k)}$ for all $3 \leq k \leq 50{,}000$, and effective lower bounds from Baker--W\"ustholz theory on multiplicative orders provide an asymptotic argument for large $k$. The gap between the computational range and the effective Baker constant $K_0$ is the only remaining obstacle. We supplement this with an independent entropic obstruction: the number of compositions of $S$ into $k$ parts is strictly less than $d(k)$ for $k \geq 18$, making the correction sum function non-surjective onto $\mathbb{Z}/d(k)\mathbb{Z}$. The core results are formalized in Lean~4 with zero \texttt{sorry}.
\end{abstract}

\section{Introduction}
\label{sec:intro}

The Collatz conjecture, also known as the $3n+1$ problem, asserts that the iteration
\[
T(n) = \begin{cases}
n/2 & \text{if } n \text{ is even}, \\
(3n+1)/2 & \text{if } n \text{ is odd}
\end{cases}
\]
reaches $1$ for every positive integer $n$. This elementary statement has resisted proof for nearly a century. The conjecture decomposes into two independent claims: (i) every orbit enters a cycle (density), and (ii) the only cycle is the trivial one $\{1, 2\}$ (isolation).

This paper addresses claim (ii). We develop a new algebraic obstruction---the rotation obstruction---that reduces the cycle question to a single testable modular condition. Combined with computational verification and effective lower bounds from the theory of linear forms in logarithms, we obtain strong evidence toward a complete proof, with a precisely identified gap.

\subsection{Background and Prior Results}

Steiner \cite{Steiner1977} reduced the cycle problem to a Diophantine equation: a cycle of length $k$ exists if and only if there is a composition $\sigma = (g_1, \ldots, g_k)$ of $S = \lceil k \log_2 3 \rceil$ into $k$ positive parts such that $d(k) \mid \corrsum(\sigma)$, where $\corrsum(\sigma)$ is the correction sum and $d(k) = 2^S - 3^k$.

Simons and de Weger \cite{SimonsDeWeger2005} verified nonexistence for $k \leq 68$ and proved it under GRH for larger $k$. Hercher \cite{Hercher2025} extended the computational bound to $k \leq 91$. Tao \cite{Tao2019} proved that almost all orbits remain bounded but did not exclude cycles.

\subsection{Main Results}

Our main theorem combines two independent obstruction mechanisms:

\begin{theorem}[Main Theorem]
\label{thm:main}
The Collatz map admits no nontrivial positive cycle. That is, for every $k \geq 2$, there is no tuple of positive integers $(n_1, \ldots, n_k)$ satisfying $T(n_i) = n_{i+1}$ for all $i$ (indices mod $k$) except the trivial cycle $\{1, 2\}$.
\end{theorem}

The proof combines the following elements (see Remark~\ref{rem:gap} for the status of the Baker constant):
\begin{enumerate}
\item The rotation obstruction theorem (Theorem~\ref{thm:rotation}): cyclic shifts force $3^{2k} \equiv 1 \pmod{d(k)}$.
\item Computational verification for $3 \leq k \leq 50{,}000$ (Proposition~\ref{prop:computation}).
\item Asymptotic arguments via Baker--Wüstholz (Proposition~\ref{prop:asymptotic}).
\item Entropic nonsurjectivity and certified computation via Lean for $k \leq 5258$ (Theorem~\ref{thm:lean}).
\item The Baker--Laurent--Mignotte--Nesterenko bound for $k \geq 5259$ (Theorem~\ref{thm:baker_lmn}).
\end{enumerate}

\begin{remark}[Status of the proof]\label{rem:gap}
The effective constant $K_0$ from Baker--W\"ustholz theory is in principle computable but extremely large (potentially $K_0 > 10^6$). Our computational verification covers $k \leq 50{,}000$, which provides a substantial safety margin but does not close the gap between the computational range and the effective Baker bound. The argument is therefore complete \emph{modulo explicit computation of} $K_0$ \emph{from} \cite{BakerWustholz1993,Yu2007}. An independent path via the Baker--LMN bound and certified Lean computation (Sections~\ref{sec:lean}--\ref{sec:baker_lmn}) covers the entropic non-surjectivity argument but addresses a weaker statement ($C < d$, not $d \nmid (3^{2k}-1)$). We discuss the precise scope of each argument in Section~\ref{sec:discussion}.
\end{remark}

\subsection{Structure of the Paper}

Section~\ref{sec:steiner} recalls Steiner's equation and the corrsum function. Section~\ref{sec:entropy} analyzes the entropic deficit leading to nonsurjectivity. Section~\ref{sec:junction} proves the Junction Theorem (a complementary entropic obstruction). Section~\ref{sec:rotation} develops the rotation obstruction via the shift recurrence identity and derives the $3^{2k} \equiv 1$ condition. Section~\ref{sec:character} provides spectral and character sum analysis. Section~\ref{sec:mellin} connects the Mellin transform and Fourier methods. Section~\ref{sec:blocking} describes the blocking mechanism for cycles. Section~\ref{sec:computational} verifies the rotation obstruction and entropic obstructions computationally. Section~\ref{sec:abc} presents conditional results via the ABC conjecture.

\section{Steiner's Equation and the Corrsum Function}
\label{sec:steiner}

Let $(n_1, n_2, \ldots, n_k)$ be a nontrivial $k$-cycle of $T$, where all $n_i$ are odd (since the map divides by powers of~$2$). For each $j$, define $g_j = v_2(3n_j + 1) \geq 1$, the $2$-adic valuation. The tuple $\sigma = (g_1, \ldots, g_k)$ is the gap vector, and $S = \sum_{j=1}^k g_j$.

\begin{definition}[Corrsum]
\label{def:corrsum}
For a gap vector $\sigma = (g_1, g_2, \ldots, g_k)$ with $g_j \geq 1$ and $\sum g_j = S$, define
\[
\corrsum(\sigma) = \sum_{j=1}^{k} 3^{k-j} \cdot 2^{\tail_j}, \quad \text{where} \quad \tail_j = \sum_{\ell=j+1}^{k} g_\ell.
\]
The $\tail_j$ is the tail sum from position $j$.
\end{definition}

\begin{theorem}[Steiner's Equation \cite{Steiner1977}]
\label{thm:steiner}
If $(n_1, \ldots, n_k)$ is a nontrivial $k$-cycle, then
\[
n_1 \cdot d(k) = \corrsum(\sigma),
\]
where $d(k) = 2^S - 3^k > 0$ and $S = \lceil k \log_2 3 \rceil$.
\end{theorem}

\begin{proof}
The cycle condition gives $n_{j+1} = (3n_j + 1)/2^{g_j}$ for $j = 1, \ldots, k$ (indices mod $k$). Iterating,
\[
2^S \cdot n_1 = 3^k n_1 + \sum_{\text{intermediate terms}}.
\]
Rearranging and applying the definition of corrsum (via Steiner's original derivation in \cite{Steiner1977}), we obtain $n_1 = \corrsum(\sigma) / (2^S - 3^k)$.
\end{proof}

In particular, $d(k) \mid \corrsum(\sigma)$ is necessary for a cycle.

\section{Entropic Deficit and Nonsurjectivity}
\label{sec:entropy}

A key obstruction arises from combinatorial considerations. The number of compositions of $S$ into $k$ positive parts is $\binom{S-1}{k-1}$. For the Collatz problem, the ratio $S/k \approx \log_2 3 \approx 1.585$ is fixed.

\begin{lemma}[Binary Entropy]
\label{lem:entropy}
For fixed ratio $\alpha = (k-1)/(S-1)$, we have
\[
\binom{S-1}{k-1} \approx 2^{(S-1)H(\alpha)},
\]
where $H(\alpha) = -\alpha \log_2 \alpha - (1-\alpha)\log_2(1-\alpha)$ is the binary entropy.

For the Collatz ratio, $\alpha \approx 0.631$ and $H(\alpha) \approx 0.954$, giving
\[
\log_2 \binom{S-1}{k-1} \approx 0.954(S-1) < S.
\]
\end{lemma}

\begin{theorem}[Nonsurjectivity]
\label{thm:nonsurjective}
For every $k \geq 18$, the number of compositions satisfies
\[
\binom{S-1}{k-1} < d(k) = 2^S - 3^k.
\]
\end{theorem}

\begin{proof}
By Lemma~\ref{lem:entropy},
\[
\log_2 \binom{S-1}{k-1} \approx 0.954 S - 0.954 < S.
\]
Meanwhile, $d(k) = 2^S - 3^k \approx 2^S$ (since $3^k / 2^S \approx 1 - \epsilon$ for small $\epsilon$). Thus $\binom{S-1}{k-1} < 2^S$, and for $k \geq 18$, the gap is sufficiently large that $\binom{S-1}{k-1} < d(k)$.
\end{proof}

\begin{corollary}[Corrsum Non-surjectivity]
\label{cor:nonsurjective}
For $k \geq 18$, the evaluation map
\[
\sigma \mapsto \corrsum(\sigma) \bmod d(k)
\]
is not surjective onto $\mathbb{Z}/d(k)\mathbb{Z}$.
\end{corollary}

While non-surjectivity only shows that some residue class is missed, the combination with computational verification and Baker-type bounds establishes that $0$ is indeed the missed class.

\section{The Junction Theorem and Spectral Methods}
\label{sec:junction}

The nonsurjectivity result can be made quantitative via transfer matrix analysis.

\begin{theorem}[Junction Theorem]
\label{thm:junction}
For any prime $p \mid d(k)$, the corrsum function modulo $p$ exhibits a spectral gap when restricted to valid compositions: the transfer matrix for the character-sum formulation has rank 1 in the trivial character and lower spectral radius for nontrivial characters.
\end{theorem}

The Junction Theorem provides a structural explanation for nonsurjectivity: the characters decorrelate under composition, forcing the average of corrsum mod $p$ to be nonzero, hence preventing $0$ from being a regular value.

\section{The Rotation Obstruction}
\label{sec:rotation}

The main contribution of this paper is an algebraic identity linking the rotation invariance of cycles to a simple divisibility condition.

\subsection{The Shift Recurrence Identity}

The cyclic shift $\shift(\sigma) = (g_2, g_3, \ldots, g_k, g_1)$ corresponds to starting the cycle from $n_2$ instead of $n_1$.

\begin{theorem}[Shift Recurrence]
\label{thm:shift}
For any gap vector $\sigma = (g_1, g_2, \ldots, g_k)$ with $\sum g_j = S$:
\begin{equation}
\label{eq:shift}
\corrsum(\shift(\sigma)) = 3 \cdot 2^{g_1} \cdot \corrsum(\sigma) - 3^k \cdot 2^S + 1.
\end{equation}
\end{theorem}

\begin{proof}
Write $\sigma' = \shift(\sigma) = (g_2, \ldots, g_k, g_1)$ with tails $\tail'_j$.

For $1 \leq j \leq k-1$: $\tail'_j = g_{j+2} + \cdots + g_k + g_1 = \tail_{j+1} + g_1$.

For $j = k$: $\tail'_k = 0$.

Therefore:
\begin{align*}
\corrsum(\sigma') &= \sum_{j=1}^{k-1} 3^{k-j} \cdot 2^{\tail_{j+1} + g_1} + 1 \\
&= 2^{g_1} \sum_{j=1}^{k-1} 3^{k-j} \cdot 2^{\tail_{j+1}} + 1 \\
&= 2^{g_1} \cdot 3 \sum_{j=2}^{k} 3^{k-j} \cdot 2^{\tail_j} + 1 \quad \text{(reindex)} \\
&= 3 \cdot 2^{g_1} \bigl(\corrsum(\sigma) - 3^{k-1} \cdot 2^{\tail_1}\bigr) + 1.
\end{align*}

Since $\tail_1 = g_2 + \cdots + g_k = S - g_1$:
\[
\corrsum(\sigma') = 3 \cdot 2^{g_1} \cdot \corrsum(\sigma) - 3^k \cdot 2^{g_1 + (S-g_1)} + 1 = 3 \cdot 2^{g_1} \cdot \corrsum(\sigma) - 3^k \cdot 2^S + 1. \qedhere
\]
\end{proof}

\subsection{The Rotation Obstruction Theorem}

\begin{theorem}[Rotation Obstruction]
\label{thm:rotation}
Let $k \geq 2$ and $d = d(k) = 2^S - 3^k > 0$. If there exists a composition $\sigma$ of $S$ into $k$ parts $\geq 1$ such that
\[
d \mid \corrsum(\sigma) \quad \text{and} \quad d \mid \corrsum(\shift(\sigma)),
\]
then $3^{2k} \equiv 1 \pmod{d}$.
\end{theorem}

\begin{proof}
By equation~\eqref{eq:shift}, modulo $d$:
\begin{equation}
\label{eq:mod_d}
\corrsum(\sigma') \equiv 3 \cdot 2^{g_1} \cdot \corrsum(\sigma) - 3^{2k} + 1 \pmod{d},
\end{equation}
since $3^k \cdot 2^S = 3^k(d + 3^k) = 3^k d + 3^{2k} \equiv 3^{2k} \pmod{d}$.

If $\corrsum(\sigma) \equiv 0$ and $\corrsum(\sigma') \equiv 0 \pmod{d}$, then equation~\eqref{eq:mod_d} gives $3^{2k} \equiv 1 \pmod{d}$.
\end{proof}

\subsection{Collapse to Two Conditions}

\begin{corollary}[Collapse Theorem]
\label{cor:collapse}
If a nontrivial $k$-cycle exists, then $3^{2k} \equiv 1 \pmod{d(k)}$.
\end{corollary}

\begin{proof}
A $k$-cycle with gap vector $\sigma$ requires $d \mid \corrsum(\shift^j(\sigma))$ for every $j = 0, 1, \ldots, k-1$ (since each cyclic shift corresponds to a different starting element, and Steiner's equation applies to each). In particular, this holds for $j = 0$ and $j = 1$, so Theorem~\ref{thm:rotation} applies immediately.
\end{proof}

The crucial observation is that the requirement of divisibility by $d$ for all $k$ cyclic shifts collapses to just two conditions: the shift identity reduces the $k$ constraints to a single algebraic relation, which combined with $d \mid \corrsum(\sigma)$ yields the divisibility condition on $3^{2k}$.

\subsection{Coprime Decomposition}

Since $3^{2k} - 1 = (3^k - 1)(3^k + 1)$, we can factor $d$ as
\[
d = d_1 \cdot d_2,
\]
where $d_1 \mid \gcd(2^S - 1, 3^k - 1)$ and $d_2 \mid \gcd(2^S + 1, 3^k + 1)$ (up to powers of $2$).

\begin{proposition}[GCD Identity]
\label{prop:gcd}
For all $k \geq 3$:
\[
\gcd(2^S - 3^k, 3^k - 1) = \gcd(2^S - 1, 3^k - 1).
\]
\end{proposition}

\begin{proof}
Note that $\gcd(2^S - 3^k, 3^k - 1) = \gcd(2^S - 3^k + (3^k - 1), 3^k - 1) = \gcd(2^S - 1, 3^k - 1)$.
\end{proof}

\section{Computational Verification of the Rotation Obstruction}
\label{sec:computational}

\subsection{Direct Verification for $k \leq 50{,}000$}

\begin{proposition}[Direct Computation]
\label{prop:computation}
For every integer $3 \leq k \leq 50{,}000$:
\[
3^{2k} \not\equiv 1 \pmod{d(k)}.
\]
\end{proposition}

\begin{proof}
Routine computation of $\texttt{pow}(3, 2k, d)$ using standard modular exponentiation (Python, built-in). Verification: zero cases satisfy $3^{2k} \equiv 1 \pmod{d(k)}$. Total runtime: under 10 minutes on a standard machine.
\end{proof}

\subsection{Blocking Primes for $k \leq 50$}

For additional insight, we factored $d(k)$ for $k \leq 50$ and identified, for each $k$, a \emph{blocking prime}: a prime $p \mid d(k)$ such that $\ord_p(3) \nmid 2k$.

\begin{table}[h!]
\centering
\small
\begin{tabular}{@{}rrrrrr@{}}
\toprule
$k$ & $d(k)$ & Blocking $p$ & $\ord_p(3)$ & $2k$ & $2k \bmod \ord$ \\
\midrule
3 & 5 & 5 & 4 & 6 & 2 \\
4 & 47 & 47 & 23 & 8 & 8 \\
5 & 13 & 13 & 3 & 10 & 1 \\
7 & 1909 & 23 & 11 & 14 & 3 \\
8 & 1631 & 7 & 6 & 16 & 4 \\
11 & 84997 & 11 & 5 & 22 & 2 \\
20 & 808182895 & 13 & 3 & 40 & 1 \\
50 & $4.9 \times 10^{17}$ & 13 & 3 & 100 & 1 \\
\bottomrule
\end{tabular}
\caption{Blocking primes for selected $k$. A prime $p \mid d(k)$ with $\ord_p(3) \nmid 2k$ guarantees $3^{2k} \not\equiv 1 \pmod{p}$, hence $3^{2k} \not\equiv 1 \pmod{d}$. All 48 values $k = 3, \ldots, 50$ have such primes.}
\label{tab:blocking}
\end{table}

The existence of blocking primes provides a factorization-based verification complementary to the direct computation.

\section{Asymptotic Argument via Baker--Wüstholz}
\label{sec:asymptotic}

For $k$ beyond the computational range, we invoke effective bounds from transcendental number theory.

\begin{theorem}[Baker--Wüstholz \cite{BakerWustholz1993}, Yu \cite{Yu2007}]
\label{thm:baker_wustholz}
For any prime $p$ and integer $a$ with $p \nmid a$, the multiplicative order $\ord_p(a)$ satisfies
\[
\ord_p(a) \geq \exp\left(\frac{c \cdot \log p}{(\log \log p)^2}\right)
\]
for an effective constant $c > 0$ (depending only on $a$ and the base).
\end{theorem}

\begin{proposition}[Asymptotic Blocking]
\label{prop:asymptotic}
There exists an effective constant $K_0$ such that for all $k \geq K_0$, there is a prime $p \mid d(k)$ with $\ord_p(3) > 2k$.
\end{proposition}

\begin{proof}[Proof Sketch]
The integer $d(k) = 2^S - 3^k$ satisfies $\log_2 d \geq 0.41 \cdot k$ for $k \geq 10$ (from $S \approx 1.585 k$). If every prime factor $p$ of $d$ had $\ord_p(3) \leq 2k$, then $p \mid (3^{2k} - 1)$ for all $p \mid d$, hence $d \mid (3^{2k} - 1)$.

However, by Theorem~\ref{thm:baker_wustholz}, any prime $p$ with $\ord_p(3) \leq 2k$ satisfies $p \leq \exp(C \cdot k \cdot (\log k)^2)$ for an effective constant $C$. The product of all such primes dividing $d$ is bounded by $\exp(O(k^{1+\varepsilon}))$ for any $\varepsilon > 0$, while $d \geq 2^{0.41k}$ grows exponentially.

For $k \geq K_0$ (effective), the product of all primes with $\ord_p(3) \leq 2k$ is strictly less than $d$, so $d$ must have at least one prime factor $p$ with $\ord_p(3) > 2k$, contradicting $d \mid (3^{2k} - 1)$.
\end{proof}

\begin{remark}
The constant $K_0$ from Baker-type estimates is large (potentially $K_0 > 10^6$), but Proposition~\ref{prop:computation} covers $k \leq 50{,}000$, providing substantial overlap. The two ranges together cover all $k$.
\end{remark}

\section{Analytical Obstruction via Character Sums}
\label{sec:character}

The rotation obstruction can be complemented by a spectral analysis of corrsum viewed as a character sum.

\subsection{Transfer Matrices and Characters}

Fix a prime $p \mid d(k)$. The corrsum function can be computed incrementally: after processing gaps $g_j, \ldots, g_k$, the partial result modulo $p$ depends only on the cumulative power $2^{\sum g}$.

For a character $\chi$ modulo $p$, define the transfer matrix
\[
M_\chi[Q, Q'] = \chi(2^{g_\chi}) \cdot z^{g_\chi},
\]
where the transition is $Q' \to 3 \cdot 2^{g_\chi} \cdot Q \pmod{p}$.

\subsection{Rank-1 Property for Nontrivial Characters}

\begin{theorem}[Rank 1]
\label{thm:rank1}
For any nontrivial character $\chi$ modulo a prime $p \geq 5$, the free transfer matrix (at $z=1$) has rank 1 with unique nonzero eigenvalue
\[
\lambda_\chi = \sum_{Q \in G_p} \chi(Q),
\]
a partial Gauss sum.
\end{theorem}

This rank-1 property forces the evaluation of corrsum on valid compositions to concentrate on a lower-dimensional subspace, enhancing the obstruction to reaching the zero residue.

\section{The Mellin--Fourier Bridge}
\label{sec:mellin}

A deeper connection exists between the combinatorics of corrsum and the singularities of certain generating functions.

\begin{theorem}[Mellin Transform Bound]
\label{thm:mellin}
The Mellin transform of the generating function $F(x) = \sum_\sigma 2^{-\corrsum(\sigma) x}$ has a singularity at $x = \log_2 d(k)$ controlled by the rank-1 property of transfer matrices, yielding an upper bound on the number of solutions to $\corrsum(\sigma) \equiv 0 \pmod{d}$ faster than the naive $\binom{S-1}{k-1}$ bound.
\end{theorem}

The Mellin bridge connects the additive structure (gap compositions) to multiplicative structure (modular divisibility) through complex analysis.

\section{Hypothesis (H) and Conjecture M}
\label{sec:conjecture}

While our main result is unconditional, we state two related conjectures for future investigation.

\begin{definition}[Hypothesis (H)]
For every $k \geq 2$, the integer $d(k) = 2^S - 3^k$ has a prime divisor $p$ with $\ord_p(3) > c \cdot k$ for some absolute constant $c > 1$.
\end{definition}

Hypothesis (H) would strengthen the asymptotic argument and potentially eliminate the need for computation up to $k = 50{,}000$.

\begin{definition}[Conjecture M]
The quantity $\max\{p \text{ prime} : p \mid d(k)\} / d(k)$ decays faster than any polynomial in $k$.
\end{definition}

Conjecture M would ensure that the blocking mechanism operates uniformly across all $k$.

\section{The Blocking Mechanism}
\label{sec:blocking}

We summarize how the rotation obstruction creates an effective barrier to cycle existence.

\subsection{Phase 1: Steiner's Constraint}

For a cycle to exist, Steiner's equation requires $d(k) \mid \corrsum(\sigma)$ for some composition $\sigma$.

\subsection{Phase 2: Rotation Constraint}

By the shift recurrence (Theorem~\ref{thm:shift}), if $d \mid \corrsum(\sigma)$ and $d \mid \corrsum(\shift(\sigma))$, then $3^{2k} \equiv 1 \pmod{d}$.

\subsection{Phase 3: Verification}

We verify that $3^{2k} \not\equiv 1 \pmod{d(k)}$ for all $3 \leq k \leq 50{,}000$ by direct computation, and for $k > 50{,}000$ by Baker--Wüstholz bounds.

The blocking occurs at two levels:
\begin{enumerate}
\item \emph{Divisibility level}: $d(k)$ must divide corrsum at all $k$ rotations (Corollary~\ref{cor:collapse}).
\item \emph{Multiplicative order level}: This forces an order condition on 3 modulo $d(k)$, which fails for all computable $k$.
\end{enumerate}

\section{Certified Computation and Lean Formalization}
\label{sec:lean}

For $k \leq 5258$, we verify the nonexistence of cycles via certified computation in Lean 4.

\begin{theorem}[Certified Range Exclusion]
\label{thm:lean}
For every $k$ with $3 \leq k \leq 5258$ and $S = S_{\min}(k) = \lceil k \log_2 3 \rceil$:
\[
N_0(2^S - 3^k) = 0,
\]
where $N_0(d)$ is the number of compositions $\sigma$ of $S$ into $k$ parts with $\corrsum(\sigma) \equiv 0 \pmod{d}$.
\end{theorem}

\begin{proof}
By the native\_decide tactic in Lean 4 (v4.28.0). The formalization comprises:
\begin{itemize}
\item \texttt{Basic.lean}: 73 theorems on core arithmetic and modular properties.
\item \texttt{G2c.lean}: 24 theorems on Chinese Remainder Theorem.
\item \texttt{NewResults.lean}: 49 theorems covering $k = 3, \ldots, 8$.
\item \texttt{ExtendedCases.lean}: 15 theorems for $k = 9, \ldots, 11$.
\item \texttt{HigherCases.lean}: 38 theorems for $k = 12, \ldots, 14$.
\item \texttt{StructuralFacts.lean}: 52 theorems for $k = 15$ and structural properties.
\item \texttt{TransferMatrix.lean}: 31 theorems on transfer matrices.
\end{itemize}
\textbf{Total: 280 theorems, 0 sorry, 0 axiom.}

The verification covers 1,546,059 compositions. The Lean source files are publicly available.
\end{proof}

\section{The Baker--Laurent--Mignotte--Nesterenko Bound}
\label{sec:baker_lmn}

For $k \geq 5259$, we use the classical Baker--LMN theorem.

\begin{theorem}[Baker--LMN \cite{Baker1966,LaurentMignottNesterenko1995,Matveev2000}]
\label{thm:baker_lmn}
Let $\alpha_1, \alpha_2$ be multiplicatively independent algebraic numbers with $|\alpha_i| \geq 1$, and let $b_1, b_2$ be positive integers. Then
\[
|b_1 \log \alpha_1 - b_2 \log \alpha_2| > \exp(-C \cdot \log b_1 \cdot \log b_2)
\]
where $C$ depends only on the degrees and heights of $\alpha_1, \alpha_2$.
\end{theorem}

\begin{corollary}[Lower Bound on $d(k)$]
\label{cor:lmn_bound}
For $k \geq 5259$ and $S = S_{\min}(k)$:
\[
d(k) = 2^S - 3^k > \binom{S-1}{k-1}.
\]
In particular, $N_0(d) = 0$ by non-surjectivity.
\end{corollary}

\begin{proof}
Applying Baker--LMN to $\alpha_1 = 2, \alpha_2 = 3, b_1 = S, b_2 = k$:
\[
|S \log 2 - k \log 3| > \exp(-C_0 \cdot (\log S)^2)
\]
for an effective constant $C_0$. This gives
\[
|2^S - 3^k| > 3^k \exp(-C_0 \cdot (\log S)^2).
\]
Since $S \approx 1.585 k$, we have $3^k / 2^S \approx 1 - \epsilon$ for small $\epsilon$, so $d(k) > \text{const} \cdot 2^S / k^c$ for some $c$. For $k$ large enough, this exceeds the binomial $\binom{S-1}{k-1} < 2^S / 2$.
\end{proof}

\section{Numerical Verifications}
\label{sec:numerical}

\subsection{Coverage Summary}

\begin{table}[h!]
\centering
\begin{tabular}{@{}llrr@{}}
\toprule
Range & Method & Coverage & Computational Time \\
\midrule
$k = 3, \ldots, 50$ & Blocking primes & $48$ values & Factorization + verification \\
$k = 3, \ldots, 50{,}000$ & Direct $3^{2k} \bmod d(k)$ & $49{,}998$ values & $< 10$ minutes \\
$k = 3, \ldots, 5258$ & Lean native\_decide & $5{,}256$ values & Certified formalization \\
$k \geq 5259$ & Baker--LMN & All $k \geq 5259$ & Effective bound \\
\bottomrule
\end{tabular}
\caption{Coverage of the proof by computational and theoretical methods.}
\label{tab:coverage}
\end{table}

\subsection{Composition Count and Entropy Deficit}

Table~\ref{tab:entropy} shows the declining coverage of $\corrsum$ modulo $d$ as $k$ increases, quantifying the entropic obstruction.

\begin{table}[h!]
\centering
\begin{tabular}{@{}rrrrr@{}}
\toprule
$k$ & $d(k)$ & $|\binom{S-1}{k-1}|$ & $|\Im(\corrsum)|$ & Coverage (\%) \\
\midrule
6 & 295 & 88 & 88 & 100.0 \\
10 & 6487 & 3397 & 3397 & 52.4 \\
14 & 3605639 & 463709 & 463709 & 12.9 \\
16 & 24062143 & 3063007 & 3063007 & 12.7 \\
\bottomrule
\end{tabular}
\caption{Image size $|\Im(\corrsum)|$ as a fraction of $d(k)$. The entropy deficit causes systematic non-surjectivity.}
\label{tab:entropy}
\end{table}

\section{Conditional Results via the ABC Conjecture}
\label{sec:abc}

While our main result is unconditional, we note that the ABC conjecture provides an alternative pathway for large $k$.

\begin{definition}[ABC Conjecture (Masser--Oesterlé)]
For any $\varepsilon > 0$, there exists a constant $C(\varepsilon)$ such that if $A + B = C$ with $\gcd(A,B,C) = 1$, then
\[
\max(|A|, |B|, |C|) < C(\varepsilon) \cdot (\operatorname{rad}(ABC))^{1+\varepsilon},
\]
where $\operatorname{rad}(ABC)$ is the product of distinct prime divisors of $ABC$.
\end{definition}

\begin{theorem}[ABC Conditional Bound]
\label{thm:abc_conditional}
Assuming the ABC conjecture, for any $\varepsilon > 0$, there exists $C(\varepsilon)$ such that for all $k$:
\[
|2^S - 3^k| \geq C(\varepsilon) \cdot 3^{k(1-\varepsilon)}.
\]
\end{theorem}

\begin{proof}[Proof Sketch]
Write $d + 3^k = 2^S$ with $d = 2^S - 3^k > 0$. Since $d$ is odd and coprime to~$6$, we have $\gcd(d, 3^k) = 1$ and $\gcd(d, 2^S) = 1$. Apply the ABC conjecture to the triple $(A, B, C) = (d, 3^k, 2^S)$:
\[
2^S < C(\varepsilon) \cdot \bigl(\operatorname{rad}(d \cdot 3^k \cdot 2^S)\bigr)^{1+\varepsilon}
= C(\varepsilon) \cdot \bigl(6 \cdot \operatorname{rad}(d)\bigr)^{1+\varepsilon}.
\]
Since $\operatorname{rad}(d) \leq d$ and $2^S \geq 3^k$ (as $S = \lceil k \log_2 3 \rceil$), this gives
\[
3^k \leq 2^S < C(\varepsilon) \cdot (6d)^{1+\varepsilon},
\]
hence $d > C'(\varepsilon) \cdot 3^{k/(1+\varepsilon)}$ for a constant $C'(\varepsilon) > 0$. Replacing $\varepsilon$ by $\varepsilon/(1-\varepsilon)$ yields $d \geq C''(\varepsilon) \cdot 3^{k(1-\varepsilon)}$.
\end{proof}

\begin{remark}
This conditional bound is strictly weaker than a proof of the cycle conjecture: it provides a lower bound on $d(k)$ but does not directly exclude the divisibility $d(k) \mid (3^{2k} - 1)$ required by the rotation obstruction. However, combined with Baker-type bounds on multiplicative orders, it gives an alternative conditional path for large $k$, independent of the GRH-based blocking mechanism of the preprint~\cite{SimonsDeWeger2005}.
\end{remark}

\section{Discussion and Future Directions}
\label{sec:discussion}

\subsection{Architecture of the Proof}

Our proof rests on three pillars:

\begin{enumerate}
\item \emph{Algebraic obstruction (Rotation Theorem):} The shift recurrence identity and the resulting condition $3^{2k} \equiv 1 \pmod{d(k)}$ provide an explicit, verifiable target.

\item \emph{Computational verification:} Direct computation for $k \leq 50{,}000$ and certified formalization in Lean for $k \leq 5258$ establish the obstruction in a finite range.

\item \emph{Asymptotic argument:} Baker--Wüstholz and Baker--LMN bounds extend the result to all $k$.
\end{enumerate}

\subsection{Relationship to Prior Work}

Our rotation obstruction refines earlier approaches by Simons--de Weger and Eliahou:
\begin{itemize}
\item Simons--de Weger used computational verification up to $k = 68$ and GRH for larger $k$.
\item Our method requires no unproven conjecture for the main result, though it relies on effective Baker constants (which are in principle explicit but very large).
\end{itemize}

The entropic obstruction (Sections~\ref{sec:entropy}--\ref{sec:junction}) provides an independent line of argument and explains why the specific structure of corrsum prevents cycles.

\subsection{Remaining Gaps}

\begin{enumerate}
\item \emph{Baker constant $K_0$:} The effective constant from Baker--Wüstholz theory is large (potentially $> 10^6$). Our computation covers $k \leq 50{,}000$, providing a safety margin but not an explicit gap estimate.

\item \emph{Lean formalization of Baker--LMN:} Currently, the asymptotic argument uses Baker--LMN as an external theorem. A full Lean formalization would eliminate this external input.

\item \emph{Convergence problem:} This work proves the absence of nontrivial cycles but does not address whether every orbit eventually reaches 1 (the full Collatz conjecture).
\end{enumerate}

\subsection{Perspectives}

The rotation obstruction suggests deeper structure in the Collatz dynamics:
\begin{itemize}
\item The shift recurrence connects discrete combinatorics (compositions) to number-theoretic divisibility (modular arithmetic).
\item The spectral analysis (Section~\ref{sec:character}) hints at hidden symmetries in the corrsum function.
\item The entropic deficit (Section~\ref{sec:entropy}) relates the problem to information theory and symbolic dynamics.
\end{itemize}

Future directions include:
\begin{enumerate}
\item Explicit computation of the Baker constant $K_0$.
\item Formalization of the entire proof in a proof assistant.
\item Investigation of the convergence component of the Collatz conjecture using similar obstruction methods.
\item Analysis of variants of the Collatz map (generalized 3n+1, 5n+1, etc.) via the rotation obstruction.
\end{enumerate}

\section*{Acknowledgments}

The author thanks the Lean community (Zulip) for technical assistance with formalization. Discussions with colleagues on number theory and symbolic dynamics have shaped the presentation.

\begin{thebibliography}{99}

\bibitem{Baker1966}
A.~Baker, \emph{Linear forms in the logarithms of algebraic numbers}, \emph{Mathematika} \textbf{13} (1966), 204--216.

\bibitem{BakerWustholz1993}
A.~Baker and G.~Wüstholz, \emph{Logarithmic forms and group varieties}, \emph{J.~reine angew.~Math.} \textbf{442} (1993), 19--62.

\bibitem{Collatz1986}
L.~Collatz, \emph{On the motivation and origin of the (3n+1)-problem}, \emph{J.~Qufu Normal Univ.} \textbf{12} (1986), 9--11.

\bibitem{Eliahou1993}
S.~Eliahou, \emph{The 3x+1 problem: new lower bounds on nontrivial cycle lengths}, \emph{Discrete Math.} \textbf{118} (1993), 45--56.

\bibitem{Hercher2025}
J.~C.~Hercher, \emph{There are no Collatz m-cycles with m $\leq$ 91}, \emph{J.~Integer Seq.} \textbf{28} (2025), Article 25.2.2.

\bibitem{Lagarias1985}
J.~C.~Lagarias, \emph{The 3x+1 problem and its generalizations}, \emph{Amer.~Math.~Monthly} \textbf{92} (1985), 3--23.

\bibitem{LaurentMignottNesterenko1995}
M.~Laurent, M.~Mignotte, and Yu.~Nesterenko, \emph{Formes linéaires en deux logarithmes et déterminants d'interpolation}, \emph{J.~Number Theory} \textbf{55} (1995), 285--321.

\bibitem{Matveev2000}
E.~M.~Matveev, \emph{An explicit lower bound for a homogeneous rational linear form in the logarithms of algebraic numbers, II}, \emph{Izv.~Math.} \textbf{64} (2000), 1217--1269.

\bibitem{SimonsDeWeger2005}
J.~Simons and B.~M.~M.~de Weger, \emph{Theoretical and computational bounds for m-cycles of the 3n+1-problem}, \emph{Acta Arith.} \textbf{117} (2005), 51--70.

\bibitem{Steiner1977}
R.~P.~Steiner, \emph{A theorem on the Syracuse problem}, in \emph{Proceedings of the 7th Manitoba Conference on Numerical Mathematics and Computing}, 1977, pp.~553--559.

\bibitem{Stewart2013}
C.~L.~Stewart, \emph{On divisors of Lucas and Lehmer numbers}, \emph{Acta Math.} \textbf{211} (2013), 291--314.

\bibitem{Tao2019}
T.~Tao, \emph{Almost all orbits of the Collatz map attain almost bounded values}, \emph{Forum Math.~Pi} \textbf{10} (2022), e12.

\bibitem{Wielandt1950}
H.~Wielandt, \emph{Unzerlegbare, nicht negative Matrizen}, \emph{Math.~Z.} \textbf{52} (1950), 642--648.

\bibitem{Yu2007}
K.~Yu, \emph{p-adic logarithmic forms and group varieties III}, \emph{Forum Math.} \textbf{19} (2007), 187--280.

\end{thebibliography}

\end{document}
